% Copyright 2009 by Tomasz Mazur
%
% This file may be distributed and/or modified in all ways.


\documentclass[xcolor=dvitex,t,11pt]{beamer}
%\documentclass[t,dvipsnames,11pt]{beamer}

% Latex Draw assets
\usepackage[usenames,dvipsnames]{pstricks}
%\usepackage{pstricks, pst-all, pstricks-add}

%%%%%%%%%%%%%%%%%%%%%%%%%%%%%%%%%%
%       SET OPTIONS BELOW        %
%%%%%%%%%%%%%%%%%%%%%%%%%%%%%%%%%%

\usetheme[
% Toggle showing page counter
pagecounter=true,
%
% String to be used between the current page and the
% total page count, e.g. of, /, from, etc.
pageofpages=of,
%
% Defines the shape of bullet points. Available options: circle, square
bullet=circle,
%
% Show a line below the frame title. 
titleline=true,
%
% Set the style of the title page (true for fancy, false for standard)
alternativetitlepage=true,
%
% Institution logo for fancy title page.
% Comment out to remove the logo from the title page.
% IMPORTANT: THERE IS A BUG IN SOME VERSIONS OF PDFLATEX AND FONTS
% ON THE LOGOS ARE NOT RENDERED PROPERLY. IN SUCH A CASE ADD `2` 
% TO THE NAME OF THE LOGO, E.G. comlab2 INSTEAD OF comlab
%titlepagelogo=images/titlepage/ou, % Phiri --October 30, 2012. Commented out
%
% Department footer logo for fancy title page
% Comment out to remove the logo from the footer of the title page/
% IMPORTANT: THERE IS A BUG IN SOME VERSIONS OF PDFLATEX AND FONTS
% ON THE LOGOS ARE NOT RENDERED PROPERLY. IN SUCH A CASE ADD `2` 
% TO THE NAME OF THE LOGO, E.G. comlab2 INSTEAD OF comlab
%titlepagefooterlogo=images/titlepage/comlab,
%
% Institution/department logo for ordinary slides
% Comment this line out to remove the logo from all the pages.
% Available logos are: ou, comlab, comlabinline, comlabou
% IMPORTANT: THERE IS A BUG IN SOME VERSIONS OF PDFLATEX AND FONTS
% ON THE LOGOS ARE NOT RENDERED PROPERLY. IN SUCH A CASE ADD `2` 
% TO THE NAME OF THE LOGO, E.G. comlab2 INSTEAD OF comlab
%ordinarypageslogo=images/comlab, % phiri --October 30, 2012. Commented out
%
%
% Add watermark in the bottom right corner
%watermark=<filename>,
%
% Set the height of the watermark.
%watermarkheight=100pt,
%
% The watermark image is 4 times bigger than watermarkheight.
%watermarkheightmult=4,
]{Torino}

% Select color theme. Available options are:
% mininmal, greenandblue, blue, red
\usecolortheme{blue}

%Select different font themes.Available options are:
% default, serif, structurebold, structureitalicserif, structuresmallcapsserif
\usefonttheme{structurebold}


%%%%%%%%%%%%%%%%%%%%%%%%%%%%%%%%%%
%       PRESENTATION INFO        %
%%%%%%%%%%%%%%%%%%%%%%%%%%%%%%%%%%


% Footer symbols
\long\def\symbolfootnote[#1]#2{%
\begingroup
\def\thefootnote{\fnsymbol{footnote}}\footnote[#1]{#2}
\endgroup
}

\author{Lighton Phiri \and Kyle Williams \and Miles Robinson\\Stuart Hammar \and Hussein Suleman}
\title{Bonolo\symbolfootnote[2, frame]{\tiny Sotho word meaning easy.}}
\subtitle{A General Digital Library System for File-based Collections}
\institute{Digital Libraries Laboratory\\Department of Computer Science\\University of Cape Town}
\date{November 13, 2012}

%
%Replace default bullets to square bullets
%
\setbeamertemplate{items}[square]

% Longtable lets you have tables that span multiple pages.
\usepackage{longtable}

%%% For tables
\usepackage{multirow}

% Rotating table headers/labels
\usepackage{rotating}

%package to enable tikZ
\usepackage{tikz}

% Footer collision with page numbers
\addtobeamertemplate{footnote}{}{\vspace{2ex}}

\usepackage{Sweave}
\begin{document}



%%%%%%%%%%%%%%%%%%%%%%%%%%%%%%%%%%
%       SLIDE DEFINITIONS        %
%%%%%%%%%%%%%%%%%%%%%%%%%%%%%%%%%%

\begin{frame}[plain]
%\frametitle{University of Cape Town}
\begin{figure}
\centering
%\framebox[\textwidth]{%
\includegraphics[width=\linewidth, height=\textwidth, keepaspectratio]{images/university-of-cape-town.eps}
%\includegraphics{images/beaultiful-universities-world.eps}
%}
%\caption{End User Interface}
\end{figure}
\end{frame}


\begin{frame}[plain]
	\titlepage
\end{frame}


\begin{frame}[fragile]
\frametitle{http://martinwest.uct.ac.za}
\begin{figure}
\centering
%\framebox[\textwidth]{%
\includegraphics[width=\linewidth, height=\textwidth, keepaspectratio]{images/icadl2012-bonolo-martinwest-collection.eps}
%}
%\caption{End User Interface}
\end{figure}
\end{frame}

\begin{frame}[fragile]
\frametitle{http://lloydbleekcollection.cs.uct.ac.za}
\begin{figure}
\centering
%\framebox[\textwidth]{%
%\includegraphics[width=\linewidth, height=\textwidth, keepaspectratio]{images/icadl2012-bonolo-lloydbleek-elephant.eps}
\includegraphics[width=\linewidth, height=\textwidth, keepaspectratio]{images/icadl2012-bonolo-bl-collection.eps}
%}
\end{figure}
\end{frame}

\begin{frame}[fragile]
\frametitle{Motivation}
\begin{itemize}
\item Preservation costs 
\begin{itemize}
\item Preservation lifecycle
\item Heritage funding model
\end{itemize}
\item Technical skills and education
\begin{itemize}
\item Content curators skillset
\item Steep learning curve for most solutions
\end{itemize}
\item Internet bandwidth
\begin{itemize}
\item Bandwidth intensive solutions
\item Cloud-centric solutions not feasible
\end{itemize}
\item Existing solutions
\begin{itemize}
\item Complexity
\end{itemize}
\end{itemize}
\end{frame}

%\begin{frame}[fragile]
%\frametitle{Previous Projects}
%\begin{itemize}
%\item Largely driven by user requirements
%\item Computing resources
%\begin{itemize}
%\item Minimal bandwidth use
%\item Minimal use of computer resources
%\item Distributable
%\end{itemize}
%\item Minimalistic approach
%\item Lightweight tools and services
%\item Projects
%\begin{itemize}
%\item Digital libraries without databases
%\item In-browser digital library services
%\item Hybrid online-offline collections
%\end{itemize}
%\end{itemize}
%\end{frame}

%\begin{frame}[fragile]
%\frametitle{Design Principles}
%\begin{itemize}
%\item Exploratory Study
%\begin{itemize}
%\item Objective of study
%\item 
%\end{itemize}
%\item Grounded Theory approach
%\begin{itemize}
%\item Data collection involved meta-analysis of existing tools
%\item Data analysis yielded eight (8) design principles
%\end{itemize}
%\item Outcome of study phase yieded eight (8) design principles
%\end{itemize}
%\end{frame}

\begin{frame}[fragile]
\frametitle{Design Principles}
\begin{itemize}
\item Design for least possible resources
\item Flexible design to facilitate extensibility
\item Hardware and/or software platform independence
\item Heterogeneous object, metadata and service integration
\item Minimalist design approach
\item Simplified preservation process
\item Structured organisation of data
\item Support for community and international standards
\end{itemize}

\end{frame}


\begin{frame}[fragile]
\frametitle{Prototype Implementation}
\begin{figure}
\centering
%\framebox[\textwidth]{%
% Generated with LaTeXDraw 2.0.8
% Sat Nov 10 04:44:28 SAST 2012
% \usepackage[usenames,dvipsnames]{pstricks}
% \usepackage{epsfig}
% \usepackage{pst-grad} % For gradients
% \usepackage{pst-plot} % For axes
\scalebox{1} % Change this value to rescale the drawing.
{
\begin{pspicture}(0,-3.1)(9.06,3.1)
\psframe[linewidth=0.04,dimen=outer](8.83,3.1)(0.83,1.9)
\psline[linewidth=0.06cm,linestyle=dashed,dash=0.16cm 0.16cm](0.63,1.6)(9.03,1.6)
\psframe[linewidth=0.04,dimen=outer](8.83,1.3)(0.83,-0.1)
\psframe[linewidth=0.03,dimen=outer](3.43,1.14)(1.03,0.66)
\psframe[linewidth=0.03,dimen=outer](6.03,1.14)(3.63,0.66)
\psframe[linewidth=0.03,dimen=outer](8.63,1.14)(6.23,0.66)
\psline[linewidth=0.06cm,linestyle=dashed,dash=0.16cm 0.16cm](0.63,-0.4)(9.03,-0.4)
%\usefont{T1}{ppl}{m}{n}
\rput(4.7876563,0.285){\scriptsize Search}
%\usefont{T1}{ppl}{m}{n}
\rput(4.785,0.885){\scriptsize Browse}
%\usefont{T1}{ppl}{m}{n}
\rput(7.398281,0.885){\scriptsize Commenting}
\psframe[linewidth=0.04,dimen=outer](8.83,-0.7)(0.83,-3.1)
\psframe[linewidth=0.03,dimen=outer,shadow=true,shadowangle=45.0,shadowsize=0.3,fillstyle=solid](3.13,-2.1)(1.53,-2.7)
%\usefont{T1}{ppl}{m}{n}
\rput(2.289375,-2.415){\scriptsize Metadata}
%\usefont{T1}{ppl}{m}{n}
\rput{-270.0}(2.646719,2.2070315){\rput(0.20078126,2.4){\small Clients}}
%\usefont{T1}{ppl}{m}{n}
\rput{-270.0}(0.69843745,0.43531248){\rput(0.1090625,0.54){\small Services}}
%\usefont{T1}{ppl}{m}{n}
\rput{-270.0}(-1.7139062,-2.0982811){\rput(0.15875,-1.9){\small Repository}}
\psframe[linewidth=0.03,dimen=outer](8.41,-1.7)(6.41,-2.1)
%\usefont{T1}{ppl}{m}{n}
\rput(7.3365626,-1.915){\scriptsize Thumbnails}
%\usefont{T1}{ppl}{m}{n}
\rput(4.7965627,-1.065){\footnotesize File System}
\psframe[linewidth=0.06,linecolor=red,linestyle=dashed,dash=0.16cm 0.16cm,dimen=outer](5.81,-1.5)(1.21,-2.9)
\psframe[linewidth=0.03,dimen=outer](8.41,-2.3)(6.41,-2.7)
%\usefont{T1}{ppl}{m}{n}
\rput(7.4223437,-2.515){\scriptsize Indicies}
\psframe[linewidth=0.06,linecolor=red,linestyle=dashed,dash=0.16cm 0.16cm,dimen=outer](8.59,-1.46)(6.19,-2.9)
\psframe[linewidth=0.03,dimen=outer](3.43,2.9)(1.43,2.1)
\psframe[linewidth=0.03,dimen=outer](5.83,2.9)(3.83,2.1)
\psframe[linewidth=0.03,linestyle=dashed,dash=0.16cm 0.16cm,dimen=outer](8.23,2.9)(6.23,2.1)
\psframe[linewidth=0.03,dimen=outer](3.43,0.54)(1.03,0.06)
\psframe[linewidth=0.03,dimen=outer](6.03,0.54)(3.63,0.06)
\psframe[linewidth=0.03,dimen=outer](8.63,0.54)(6.23,0.06)
%\usefont{T1}{ppl}{m}{n}
\rput(2.2157812,0.885){\scriptsize Authentication}
%\usefont{T1}{ppl}{m}{n}
\rput(2.1823437,0.285){\scriptsize Indexing}
%\usefont{T1}{ppl}{m}{n}
\rput(7.3765626,0.285){\scriptsize Tagging}
%\usefont{T1}{ppl}{m}{n}
\rput(2.4167187,2.495){\footnotesize Curator}
%\usefont{T1}{ppl}{m}{n}
\rput(4.83,2.495){\footnotesize End User}
%\usefont{T1}{ppl}{m}{n}
\rput(7.235156,2.495){\footnotesize Solr}
\psframe[linewidth=0.03,dimen=outer,shadow=true,shadowangle=45.0,shadowsize=0.3,fillstyle=solid](5.33,-2.1)(3.73,-2.7)
%\usefont{T1}{ppl}{m}{n}
\rput(4.54375,-2.415){\scriptsize Content}
\end{pspicture} 
}


%}
\end{figure}
\end{frame}

\begin{frame}[fragile]
\frametitle{Tools and Technologies}
\begin{itemize}
\item Core technologies
\begin{itemize}
\item Apache Tomcat
\item CSS
\item HTML
\item Java
\item JavaScript
\item XML
\end{itemize}
\item Additional software components
\begin{itemize}
\item Apache Solr
\item ImageMagick
\item SQLite
\end{itemize}
\end{itemize}
\end{frame}


\begin{frame}[fragile]
\frametitle{Repository Design}
\begin{figure}
\centering
%\framebox[\textwidth]{%
% Generated with LaTeXDraw 2.0.8
% Sun Nov 04 22:48:23 SAST 2012
% \usepackage[usenames,dvipsnames]{pstricks}
% \usepackage{epsfig}
% \usepackage{pst-grad} % For gradients
% \usepackage{pst-plot} % For axes
\scalebox{1} % Change this value to rescale the drawing.
{
\begin{pspicture}(0,-0.92)(8.0,0.92)
\psframe[linewidth=0.04,dimen=outer](8.0,0.92)(0.0,-0.92)
\psframe[linewidth=0.03,dimen=outer,shadow=true,shadowangle=45.0,shadowsize=0.3,fillstyle=solid](2.2,0.06)(0.6,-0.54)
%\usefont{T1}{ppl}{m}{n}
\rput(1.399375,-0.255){\scriptsize Metadata}
\psframe[linewidth=0.03,dimen=outer](7.58,0.46)(5.58,0.06)
%\usefont{T1}{ppl}{m}{n}
\rput(6.5265627,0.245){\scriptsize Thumbnails}
\psframe[linewidth=0.06,linecolor=red,linestyle=dashed,dash=0.16cm 0.16cm,dimen=outer](4.98,0.66)(0.38,-0.74)
\psframe[linewidth=0.03,dimen=outer](7.58,-0.14)(5.58,-0.54)
%\usefont{T1}{ppl}{m}{n}
\rput(6.592344,-0.355){\scriptsize Indicies}
\psframe[linewidth=0.06,linestyle=dashed,dash=0.16cm 0.16cm,dimen=outer](7.76,0.66)(5.36,-0.74)
\psframe[linewidth=0.03,dimen=outer,shadow=true,shadowangle=45.0,shadowsize=0.3,fillstyle=solid](4.6,0.06)(3.0,-0.54)
%\usefont{T1}{ppl}{m}{n}
\rput(3.79375,-0.255){\scriptsize Content}
\end{pspicture} 
}


%}
\end{figure}
\begin{itemize}
\item Hierarchical structure
\item Digital content and metadata stored on file system
\item Objects stored with corresponding metadata record
\item Metadata records encoded using XML
\end{itemize}
\begin{figure}
%\includegraphics[height=100pt]{images/titlepage/ou}
%\caption{Oxford University logo}
\end{figure}
\end{frame}

\begin{frame}[fragile]
\frametitle{Repository Design (2)}
\begin{figure}
\centering
%\framebox[\textwidth]{%
% Generated with LaTeXDraw 2.0.8
% Mon Nov 12 23:19:59 CST 2012
% \usepackage[usenames,dvipsnames]{pstricks}
% \usepackage{epsfig}
% \usepackage{pst-grad} % For gradients
% \usepackage{pst-plot} % For axes
\scalebox{1} % Change this value to rescale the drawing.
{
\begin{pspicture}(0,-3.0332813)(6.75,3.0332813)
\psline[linewidth=0.03cm](0.0484375,2.6710937)(0.0484375,-2.9289062)
%\usefont{T1}{ptm}{m}{n}
\rput(0.4575,2.8710938){\footnotesize archive/}
\psline[linewidth=0.03cm](0.0484375,2.3710938)(0.4484375,2.3710938)
%\usefont{T1}{ptm}{m}{n}
\rput(1.2842188,2.3910937){\footnotesize notebooks}
\psline[linewidth=0.03cm](0.7484375,2.0710938)(0.7484375,-1.9289062)
\psline[linewidth=0.03cm](0.7484375,1.8710938)(1.1484375,1.8710938)
%\usefont{T1}{ptm}{m}{n}
\rput(2.7640624,1.8910937){\footnotesize wilhelmbleeknotebooks}
\psline[linewidth=0.03cm](1.4484375,1.5710938)(1.4484375,-1.4289062)
\psline[linewidth=0.03cm](1.4484375,1.3710938)(1.8484375,1.3710938)
%\usefont{T1}{ptm}{m}{n}
\rput(3.30625,1.3910937){\footnotesize BC\_151\_A1\_4\_001}
\psline[linewidth=0.03cm](2.1484375,1.0710938)(2.1484375,-0.62890625)
\psline[linewidth=0.03cm](2.1484375,0.87109375)(2.5484376,0.87109375)
%\usefont{T1}{ptm}{m}{n}
\rput(4.0404687,0.89109373){\footnotesize A1\_4\_1\_00001.JPG}
\psline[linewidth=0.03cm](2.1484375,0.37109375)(2.5484376,0.37109375)
%\usefont{T1}{ptm}{m}{n}
\rput(4.7279687,0.39109376){\footnotesize \psframebox[linewidth=0.04,linecolor=red,linestyle=dashed,dash=0.16cm 0.16cm,framesep=0.06]{A1\_4\_1\_00001.JPG.metadata}}
\psline[linewidth=0.03cm](2.1484375,-0.12890625)(2.5484376,-0.12890625)
%\usefont{T1}{ptm}{m}{n}
\rput(2.768125,-0.10390625){:}
\psline[linewidth=0.03cm](1.4484375,-0.92890626)(1.8484375,-0.92890626)
%\usefont{T1}{ptm}{m}{n}
\rput(3.9165626,-0.92890626){\footnotesize \psframebox[linewidth=0.04,linecolor=red,linestyle=dashed,dash=0.16cm 0.16cm,framesep=0.06]{BC\_151\_A1\_4\_001.metadata}}
\psline[linewidth=0.03cm](1.4484375,-1.4289062)(1.8484375,-1.4289062)
%\usefont{T1}{ptm}{m}{n}
\rput(2.068125,-1.4239062){:}
\psline[linewidth=0.03cm](0.7484375,-1.9289062)(1.1484375,-1.9289062)
%\usefont{T1}{ptm}{m}{n}
\rput(3.458125,-1.9089062){\footnotesize \psframebox[linewidth=0.04,linecolor=red,linestyle=dashed,dash=0.16cm 0.16cm,framesep=0.06]{wilhelmbleeknotebooks.metadata}}
\psline[linewidth=0.03cm](0.0484375,-2.5289063)(0.4484375,-2.5289063)
%\usefont{T1}{ptm}{m}{n}
\rput(1.1798438,-2.5089064){\footnotesize drawings}
\psline[linewidth=0.03cm](0.0484375,-2.9289062)(0.4484375,-2.9289062)
%\usefont{T1}{ptm}{m}{n}
\rput(1.0273438,-2.8889062){\footnotesize stories}
\end{pspicture} 
}


%}
\end{figure}
\end{frame}

\begin{frame}[fragile]
\frametitle{End User Interface}
\begin{figure}
\centering
%\framebox[\textwidth]{%
% Generated with LaTeXDraw 2.0.8
% Fri Nov 09 12:53:34 SAST 2012
% \usepackage[usenames,dvipsnames]{pstricks}
% \usepackage{epsfig}
% \usepackage{pst-grad} % For gradients
% \usepackage{pst-plot} % For axes
\scalebox{1} % Change this value to rescale the drawing.
{
\begin{pspicture}(0,-0.6)(8.0,0.6)
\psframe[linewidth=0.04,dimen=outer](8.0,0.6)(0.0,-0.6)
\psframe[linewidth=0.03,dimen=outer](2.6,0.4)(0.6,-0.4)
\psframe[linewidth=0.06,linecolor=red,linestyle=dashed,dash=0.16cm 0.16cm,dimen=outer](5.0,0.4)(3.0,-0.4)
\psframe[linewidth=0.03,dimen=outer](7.4,0.4)(5.4,-0.4)
%\usefont{T1}{ppl}{m}{n}
\rput(1.5867188,-0.0050){\footnotesize Curator}
%\usefont{T1}{ppl}{m}{n}
\rput(4.0,-0.0050){\footnotesize End User}
%\usefont{T1}{ppl}{m}{n}
\rput(6.405156,-0.0050){\footnotesize Solr}
\end{pspicture} 
}


%}
%\caption{End User Interface}
\end{figure}
\begin{itemize}
\item Authentication
\item Browse
\item Commenting
\item Search
\item Tagging
\end{itemize}
\end{frame}

\begin{frame}[fragile]
\frametitle{End User Interface (2)}
\begin{figure}
\centering
%\framebox[\textwidth]{%
\includegraphics[width=\linewidth, height=\textwidth, keepaspectratio]{images/icadl2012-bonolo-ui-item.eps}
%}
%\caption{End User Interface}
\end{figure}
\end{frame}

\begin{frame}[fragile]
\frametitle{Curator Interface}
\begin{figure}
\centering
%\framebox[\textwidth]{%
% Generated with LaTeXDraw 2.0.8
% Fri Nov 09 13:07:32 SAST 2012
% \usepackage[usenames,dvipsnames]{pstricks}
% \usepackage{epsfig}
% \usepackage{pst-grad} % For gradients
% \usepackage{pst-plot} % For axes
\scalebox{1} % Change this value to rescale the drawing.
{
\begin{pspicture}(0,-0.6)(8.0,0.6)
\psframe[linewidth=0.04,dimen=outer](8.0,0.6)(0.0,-0.6)
\psframe[linewidth=0.06,linecolor=red,linestyle=dashed,dash=0.16cm 0.16cm,dimen=outer](2.6,0.4)(0.6,-0.4)
\psframe[linewidth=0.03,dimen=outer](5.0,0.4)(3.0,-0.4)
\psframe[linewidth=0.06,linecolor=red,linestyle=dashed,dash=0.16cm 0.16cm,dimen=outer](7.4,0.4)(5.4,-0.4)
%\usefont{T1}{ppl}{m}{n}
\rput(1.5867188,-0.0050){\footnotesize Curator}
%\usefont{T1}{ppl}{m}{n}
\rput(4.0,-0.0050){\footnotesize End User}
%\usefont{T1}{ppl}{m}{n}
\rput(6.405156,-0.0050){\footnotesize Solr}
\end{pspicture} 
}


%}
%\caption{End User Interface}
\end{figure}
\begin{itemize}
\item Authentication
\item Batch import and export
\item Browse
\item Collection and Object management (add, edit, delete)
\item Indexing
\item Search
\end{itemize}
\end{frame}

\begin{frame}[fragile]
\frametitle{Curator Interface (2)}
\begin{figure}
\centering
%\framebox[\textwidth]{%
\includegraphics[width=\linewidth, height=\textwidth, keepaspectratio]{images/icadl2012-bonolo-ci-item.eps}
%}
%\caption{Curator Interface}
\end{figure}
\end{frame}

\begin{frame}[fragile]
\frametitle{End User Interface UX Experiment}
\begin{itemize}
\item Objective
\begin{itemize}
\item Feasibility of building user interface using file store
\item Assess impact of file store on user experience
\end{itemize}
\item Target Group
\begin{itemize}
\item Individuals with basic Web experience
\item 17 participants (80\% students)
\end{itemize}
\item Approach
\begin{itemize}
%\item Pilot-test
\item Informed consent
\item Pre-experiment briefing
\item Typical DL end user tasks (search, browse)
\item Post-experiment questionnaire
\end{itemize}
\end{itemize}
\end{frame}


\begin{frame}[fragile]
\frametitle{End User Interface UX Experiment (2)}
\centering
% Created by tikzDevice version 0.6.2-92-0ad2792 on 2012-11-10 03:17:29
% !TEX encoding = UTF-8 Unicode
\begin{tikzpicture}[x=1pt,y=1pt]
\definecolor[named]{fillColor}{rgb}{1.00,1.00,1.00}
\path[use as bounding box,fill=fillColor,fill opacity=0.00] (0,0) rectangle (332.44,195.13);
\begin{scope}
\path[clip] (  0.00,  0.00) rectangle (332.44,195.13);
\definecolor[named]{drawColor}{rgb}{0.00,0.00,0.00}
\definecolor[named]{fillColor}{rgb}{1.00,0.00,0.00}

\path[draw=drawColor,line width= 0.4pt,line join=round,line cap=round,fill=fillColor] ( 38.71, 28.80) rectangle ( 45.15,156.10);
\definecolor[named]{fillColor}{rgb}{1.00,0.86,0.00}

\path[draw=drawColor,line width= 0.4pt,line join=round,line cap=round,fill=fillColor] ( 45.15, 28.80) rectangle ( 51.59,175.20);
\definecolor[named]{fillColor}{rgb}{0.29,1.00,0.00}

\path[draw=drawColor,line width= 0.4pt,line join=round,line cap=round,fill=fillColor] ( 51.59, 28.80) rectangle ( 58.02,172.02);
\definecolor[named]{fillColor}{rgb}{0.00,1.00,0.57}

\path[draw=drawColor,line width= 0.4pt,line join=round,line cap=round,fill=fillColor] ( 58.02, 28.80) rectangle ( 64.46,178.38);
\definecolor[named]{fillColor}{rgb}{0.00,0.57,1.00}

\path[draw=drawColor,line width= 0.4pt,line join=round,line cap=round,fill=fillColor] ( 64.46, 28.80) rectangle ( 70.90,184.75);
\definecolor[named]{fillColor}{rgb}{0.29,0.00,1.00}

\path[draw=drawColor,line width= 0.4pt,line join=round,line cap=round,fill=fillColor] ( 70.90, 28.80) rectangle ( 77.33,178.38);
\definecolor[named]{fillColor}{rgb}{1.00,0.00,0.86}

\path[draw=drawColor,line width= 0.4pt,line join=round,line cap=round,fill=fillColor] ( 77.33, 28.80) rectangle ( 83.77,181.56);
\definecolor[named]{fillColor}{rgb}{1.00,0.00,0.00}

\path[draw=drawColor,line width= 0.4pt,line join=round,line cap=round,fill=fillColor] (106.30, 28.80) rectangle (112.74, 28.80);
\definecolor[named]{fillColor}{rgb}{1.00,0.86,0.00}

\path[draw=drawColor,line width= 0.4pt,line join=round,line cap=round,fill=fillColor] (112.74, 28.80) rectangle (119.17,168.83);
\definecolor[named]{fillColor}{rgb}{0.29,1.00,0.00}

\path[draw=drawColor,line width= 0.4pt,line join=round,line cap=round,fill=fillColor] (119.17, 28.80) rectangle (125.61,168.83);
\definecolor[named]{fillColor}{rgb}{0.00,1.00,0.57}

\path[draw=drawColor,line width= 0.4pt,line join=round,line cap=round,fill=fillColor] (125.61, 28.80) rectangle (132.05,172.02);
\definecolor[named]{fillColor}{rgb}{0.00,0.57,1.00}

\path[draw=drawColor,line width= 0.4pt,line join=round,line cap=round,fill=fillColor] (132.05, 28.80) rectangle (138.48, 28.80);
\definecolor[named]{fillColor}{rgb}{0.29,0.00,1.00}

\path[draw=drawColor,line width= 0.4pt,line join=round,line cap=round,fill=fillColor] (138.48, 28.80) rectangle (144.92, 28.80);
\definecolor[named]{fillColor}{rgb}{1.00,0.00,0.86}

\path[draw=drawColor,line width= 0.4pt,line join=round,line cap=round,fill=fillColor] (144.92, 28.80) rectangle (151.36, 28.80);
\definecolor[named]{fillColor}{rgb}{1.00,0.00,0.00}

\path[draw=drawColor,line width= 0.4pt,line join=round,line cap=round,fill=fillColor] (173.89, 28.80) rectangle (180.32,168.83);
\definecolor[named]{fillColor}{rgb}{1.00,0.86,0.00}

\path[draw=drawColor,line width= 0.4pt,line join=round,line cap=round,fill=fillColor] (180.32, 28.80) rectangle (186.76,168.83);
\definecolor[named]{fillColor}{rgb}{0.29,1.00,0.00}

\path[draw=drawColor,line width= 0.4pt,line join=round,line cap=round,fill=fillColor] (186.76, 28.80) rectangle (193.20, 28.80);
\definecolor[named]{fillColor}{rgb}{0.00,1.00,0.57}

\path[draw=drawColor,line width= 0.4pt,line join=round,line cap=round,fill=fillColor] (193.20, 28.80) rectangle (199.63,172.02);
\definecolor[named]{fillColor}{rgb}{0.00,0.57,1.00}

\path[draw=drawColor,line width= 0.4pt,line join=round,line cap=round,fill=fillColor] (199.63, 28.80) rectangle (206.07,181.56);
\definecolor[named]{fillColor}{rgb}{0.29,0.00,1.00}

\path[draw=drawColor,line width= 0.4pt,line join=round,line cap=round,fill=fillColor] (206.07, 28.80) rectangle (212.51, 28.80);
\definecolor[named]{fillColor}{rgb}{1.00,0.00,0.86}

\path[draw=drawColor,line width= 0.4pt,line join=round,line cap=round,fill=fillColor] (212.51, 28.80) rectangle (218.94, 28.80);
\definecolor[named]{fillColor}{rgb}{1.00,0.00,0.00}

\path[draw=drawColor,line width= 0.4pt,line join=round,line cap=round,fill=fillColor] (241.47, 28.80) rectangle (247.91,172.02);
\definecolor[named]{fillColor}{rgb}{1.00,0.86,0.00}

\path[draw=drawColor,line width= 0.4pt,line join=round,line cap=round,fill=fillColor] (247.91, 28.80) rectangle (254.35,168.83);
\definecolor[named]{fillColor}{rgb}{0.29,1.00,0.00}

\path[draw=drawColor,line width= 0.4pt,line join=round,line cap=round,fill=fillColor] (254.35, 28.80) rectangle (260.78,172.02);
\definecolor[named]{fillColor}{rgb}{0.00,1.00,0.57}

\path[draw=drawColor,line width= 0.4pt,line join=round,line cap=round,fill=fillColor] (260.78, 28.80) rectangle (267.22, 28.80);
\definecolor[named]{fillColor}{rgb}{0.00,0.57,1.00}

\path[draw=drawColor,line width= 0.4pt,line join=round,line cap=round,fill=fillColor] (267.22, 28.80) rectangle (273.66, 28.80);
\definecolor[named]{fillColor}{rgb}{0.29,0.00,1.00}

\path[draw=drawColor,line width= 0.4pt,line join=round,line cap=round,fill=fillColor] (273.66, 28.80) rectangle (280.09, 28.80);
\definecolor[named]{fillColor}{rgb}{1.00,0.00,0.86}

\path[draw=drawColor,line width= 0.4pt,line join=round,line cap=round,fill=fillColor] (280.09, 28.80) rectangle (286.53, 28.80);
\end{scope}
\begin{scope}
\path[clip] (  0.00,  0.00) rectangle (332.44,195.13);
\definecolor[named]{drawColor}{rgb}{0.00,0.00,0.00}

\node[text=drawColor,anchor=base,inner sep=0pt, outer sep=0pt, scale=  0.60] at ( 61.24, 15.84) {Search\& Browse};

\node[text=drawColor,anchor=base,inner sep=0pt, outer sep=0pt, scale=  0.60] at (128.83, 15.84) {Navigability};

\node[text=drawColor,anchor=base,inner sep=0pt, outer sep=0pt, scale=  0.60] at (196.41, 15.84) {Similarity};

\node[text=drawColor,anchor=base,inner sep=0pt, outer sep=0pt, scale=  0.60] at (264.00, 15.84) {General};
\end{scope}
\begin{scope}
\path[clip] (  0.00,  0.00) rectangle (332.44,195.13);
\definecolor[named]{drawColor}{rgb}{0.00,0.00,0.00}

\node[text=drawColor,anchor=base,inner sep=0pt, outer sep=0pt, scale=  0.60] at (162.62,  1.44) {Average UX Ratings};

\node[text=drawColor,rotate= 90.00,anchor=base,inner sep=0pt, outer sep=0pt, scale=  0.60] at (  5.76,108.36) {Survey Participants};
\end{scope}
\begin{scope}
\path[clip] (  0.00,  0.00) rectangle (332.44,195.13);
\definecolor[named]{drawColor}{rgb}{0.00,0.00,0.00}

\path[draw=drawColor,line width= 0.4pt,line join=round,line cap=round] ( 28.80, 28.80) -- ( 28.80,187.93);

\path[draw=drawColor,line width= 0.4pt,line join=round,line cap=round] ( 28.80, 28.80) -- ( 25.20, 28.80);

\path[draw=drawColor,line width= 0.4pt,line join=round,line cap=round] ( 28.80, 60.63) -- ( 25.20, 60.63);

\path[draw=drawColor,line width= 0.4pt,line join=round,line cap=round] ( 28.80, 92.45) -- ( 25.20, 92.45);

\path[draw=drawColor,line width= 0.4pt,line join=round,line cap=round] ( 28.80,124.28) -- ( 25.20,124.28);

\path[draw=drawColor,line width= 0.4pt,line join=round,line cap=round] ( 28.80,156.10) -- ( 25.20,156.10);

\path[draw=drawColor,line width= 0.4pt,line join=round,line cap=round] ( 28.80,187.93) -- ( 25.20,187.93);

\node[text=drawColor,anchor=base east,inner sep=0pt, outer sep=0pt, scale=  0.60] at ( 21.60, 26.73) {0};

\node[text=drawColor,anchor=base east,inner sep=0pt, outer sep=0pt, scale=  0.60] at ( 21.60, 58.56) {1};

\node[text=drawColor,anchor=base east,inner sep=0pt, outer sep=0pt, scale=  0.60] at ( 21.60, 90.39) {2};

\node[text=drawColor,anchor=base east,inner sep=0pt, outer sep=0pt, scale=  0.60] at ( 21.60,122.21) {3};

\node[text=drawColor,anchor=base east,inner sep=0pt, outer sep=0pt, scale=  0.60] at ( 21.60,154.04) {4};

\node[text=drawColor,anchor=base east,inner sep=0pt, outer sep=0pt, scale=  0.60] at ( 21.60,185.86) {5};
\end{scope}
\begin{scope}
\path[clip] (  0.00,  0.00) rectangle (332.44,195.13);
\definecolor[named]{drawColor}{rgb}{0.00,0.00,0.00}

\path[draw=drawColor,line width= 0.4pt,line join=round,line cap=round] (280.09,187.93) rectangle (323.16,130.33);
\definecolor[named]{fillColor}{rgb}{1.00,0.00,0.00}

\path[draw=drawColor,line width= 0.4pt,line join=round,line cap=round,fill=fillColor] (285.49,182.53) rectangle (289.81,178.93);
\definecolor[named]{fillColor}{rgb}{1.00,0.86,0.00}

\path[draw=drawColor,line width= 0.4pt,line join=round,line cap=round,fill=fillColor] (285.49,175.33) rectangle (289.81,171.73);
\definecolor[named]{fillColor}{rgb}{0.29,1.00,0.00}

\path[draw=drawColor,line width= 0.4pt,line join=round,line cap=round,fill=fillColor] (285.49,168.13) rectangle (289.81,164.53);
\definecolor[named]{fillColor}{rgb}{0.00,1.00,0.57}

\path[draw=drawColor,line width= 0.4pt,line join=round,line cap=round,fill=fillColor] (285.49,160.93) rectangle (289.81,157.33);
\definecolor[named]{fillColor}{rgb}{0.00,0.57,1.00}

\path[draw=drawColor,line width= 0.4pt,line join=round,line cap=round,fill=fillColor] (285.49,153.73) rectangle (289.81,150.13);
\definecolor[named]{fillColor}{rgb}{0.29,0.00,1.00}

\path[draw=drawColor,line width= 0.4pt,line join=round,line cap=round,fill=fillColor] (285.49,146.53) rectangle (289.81,142.93);
\definecolor[named]{fillColor}{rgb}{1.00,0.00,0.86}

\path[draw=drawColor,line width= 0.4pt,line join=round,line cap=round,fill=fillColor] (285.49,139.33) rectangle (289.81,135.73);

\node[text=drawColor,anchor=base west,inner sep=0pt, outer sep=0pt, scale=  0.60] at (295.21,178.66) {Intuitive};

\node[text=drawColor,anchor=base west,inner sep=0pt, outer sep=0pt, scale=  0.60] at (295.21,171.46) {Simple};

\node[text=drawColor,anchor=base west,inner sep=0pt, outer sep=0pt, scale=  0.60] at (295.21,164.26) {Satisfied};

\node[text=drawColor,anchor=base west,inner sep=0pt, outer sep=0pt, scale=  0.60] at (295.21,157.06) {Expected};

\node[text=drawColor,anchor=base west,inner sep=0pt, outer sep=0pt, scale=  0.60] at (295.21,149.86) {Response};

\node[text=drawColor,anchor=base west,inner sep=0pt, outer sep=0pt, scale=  0.60] at (295.21,142.66) {Effective};

\node[text=drawColor,anchor=base west,inner sep=0pt, outer sep=0pt, scale=  0.60] at (295.21,135.46) {Beneficial};
\end{scope}
\end{tikzpicture}

\end{frame}

\begin{frame}[fragile]
\frametitle{Curator Interface UX Experiment}
\begin{itemize}
\item Objective
\begin{itemize}
\item Assess user experience when performing curation tasks
\end{itemize}
\item Target Group
\begin{itemize}
\item Individuals with no experience working with DL tools
\item Social networking site recruitment
\item 23 participants
\end{itemize}
\item Approach
\begin{itemize}
%\item Pilot-test
\item Intrinsic Motivation Inventory
\item Five (5) minute 'HOWTO' screencast
\item Curation tasks with two datasets
\item Online questionnaire
\end{itemize}
\end{itemize}
\end{frame}


\begin{frame}[fragile]
\frametitle{Curator Interface UX Experiment (2)}
\centering
% Created by tikzDevice version 0.6.2-92-0ad2792 on 2012-11-10 04:50:24
% !TEX encoding = UTF-8 Unicode
\begin{tikzpicture}[x=1pt,y=1pt]
\definecolor[named]{fillColor}{rgb}{1.00,1.00,1.00}
\path[use as bounding box,fill=fillColor,fill opacity=0.00] (0,0) rectangle (332.44,195.13);
\begin{scope}
\path[clip] (  0.00,  0.00) rectangle (332.44,195.13);
\definecolor[named]{drawColor}{rgb}{0.00,0.00,0.00}
\definecolor[named]{fillColor}{rgb}{1.00,0.00,0.00}

\path[draw=drawColor,line width= 0.4pt,line join=round,line cap=round,fill=fillColor] ( 36.85, 28.80) rectangle ( 96.01, 41.53);
\definecolor[named]{fillColor}{rgb}{0.80,1.00,0.00}

\path[draw=drawColor,line width= 0.4pt,line join=round,line cap=round,fill=fillColor] ( 36.85, 41.53) rectangle ( 96.01, 66.99);
\definecolor[named]{fillColor}{rgb}{0.00,1.00,0.40}

\path[draw=drawColor,line width= 0.4pt,line join=round,line cap=round,fill=fillColor] ( 36.85, 66.99) rectangle ( 96.01,137.01);
\definecolor[named]{fillColor}{rgb}{0.00,0.40,1.00}

\path[draw=drawColor,line width= 0.4pt,line join=round,line cap=round,fill=fillColor] ( 36.85,137.01) rectangle ( 96.01,162.47);
\definecolor[named]{fillColor}{rgb}{0.80,0.00,1.00}

\path[draw=drawColor,line width= 0.4pt,line join=round,line cap=round,fill=fillColor] ( 36.85,162.47) rectangle ( 96.01,168.83);
\definecolor[named]{fillColor}{rgb}{1.00,0.00,0.00}

\path[draw=drawColor,line width= 0.4pt,line join=round,line cap=round,fill=fillColor] (107.84, 28.80) rectangle (167.00, 47.90);
\definecolor[named]{fillColor}{rgb}{0.80,1.00,0.00}

\path[draw=drawColor,line width= 0.4pt,line join=round,line cap=round,fill=fillColor] (107.84, 47.90) rectangle (167.00,143.37);
\definecolor[named]{fillColor}{rgb}{0.00,1.00,0.40}

\path[draw=drawColor,line width= 0.4pt,line join=round,line cap=round,fill=fillColor] (107.84,143.37) rectangle (167.00,156.10);
\definecolor[named]{fillColor}{rgb}{0.00,0.40,1.00}

\path[draw=drawColor,line width= 0.4pt,line join=round,line cap=round,fill=fillColor] (107.84,156.10) rectangle (167.00,162.47);
\definecolor[named]{fillColor}{rgb}{0.80,0.00,1.00}

\path[draw=drawColor,line width= 0.4pt,line join=round,line cap=round,fill=fillColor] (107.84,162.47) rectangle (167.00,168.83);
\definecolor[named]{fillColor}{rgb}{1.00,0.00,0.00}

\path[draw=drawColor,line width= 0.4pt,line join=round,line cap=round,fill=fillColor] (178.83, 28.80) rectangle (238.00, 79.72);
\definecolor[named]{fillColor}{rgb}{0.80,1.00,0.00}

\path[draw=drawColor,line width= 0.4pt,line join=round,line cap=round,fill=fillColor] (178.83, 79.72) rectangle (238.00,149.74);
\definecolor[named]{fillColor}{rgb}{0.00,1.00,0.40}

\path[draw=drawColor,line width= 0.4pt,line join=round,line cap=round,fill=fillColor] (178.83,149.74) rectangle (238.00,156.10);
\definecolor[named]{fillColor}{rgb}{0.00,0.40,1.00}

\path[draw=drawColor,line width= 0.4pt,line join=round,line cap=round,fill=fillColor] (178.83,156.10) rectangle (238.00,168.83);
\definecolor[named]{fillColor}{rgb}{0.80,0.00,1.00}

\path[draw=drawColor,line width= 0.4pt,line join=round,line cap=round,fill=fillColor] (178.83,168.83) rectangle (238.00,168.83);
\end{scope}
\begin{scope}
\path[clip] (  0.00,  0.00) rectangle (332.44,195.13);
\definecolor[named]{drawColor}{rgb}{0.00,0.00,0.00}

\node[text=drawColor,anchor=base,inner sep=0pt, outer sep=0pt, scale=  0.60] at ( 66.43, 15.84) {Interest/Enjoyment};

\node[text=drawColor,anchor=base,inner sep=0pt, outer sep=0pt, scale=  0.60] at (137.42, 15.84) {Perceived Competence};

\node[text=drawColor,anchor=base,inner sep=0pt, outer sep=0pt, scale=  0.60] at (208.42, 15.84) {Value/Usefulness};
\end{scope}
\begin{scope}
\path[clip] (  0.00,  0.00) rectangle (332.44,195.13);
\definecolor[named]{drawColor}{rgb}{0.00,0.00,0.00}

\node[text=drawColor,anchor=base,inner sep=0pt, outer sep=0pt, scale=  0.60] at (137.42,  1.44) {Average UX Ratings};

\node[text=drawColor,rotate= 90.00,anchor=base,inner sep=0pt, outer sep=0pt, scale=  0.60] at (  5.76,108.36) {Survey Participants};
\end{scope}
\begin{scope}
\path[clip] (  0.00,  0.00) rectangle (332.44,195.13);
\definecolor[named]{drawColor}{rgb}{0.00,0.00,0.00}

\path[draw=drawColor,line width= 0.4pt,line join=round,line cap=round] ( 28.80, 28.80) -- ( 28.80,187.93);

\path[draw=drawColor,line width= 0.4pt,line join=round,line cap=round] ( 28.80, 28.80) -- ( 25.20, 28.80);

\path[draw=drawColor,line width= 0.4pt,line join=round,line cap=round] ( 28.80, 60.63) -- ( 25.20, 60.63);

\path[draw=drawColor,line width= 0.4pt,line join=round,line cap=round] ( 28.80, 92.45) -- ( 25.20, 92.45);

\path[draw=drawColor,line width= 0.4pt,line join=round,line cap=round] ( 28.80,124.28) -- ( 25.20,124.28);

\path[draw=drawColor,line width= 0.4pt,line join=round,line cap=round] ( 28.80,156.10) -- ( 25.20,156.10);

\path[draw=drawColor,line width= 0.4pt,line join=round,line cap=round] ( 28.80,187.93) -- ( 25.20,187.93);

\node[text=drawColor,anchor=base east,inner sep=0pt, outer sep=0pt, scale=  0.60] at ( 21.60, 26.73) {0};

\node[text=drawColor,anchor=base east,inner sep=0pt, outer sep=0pt, scale=  0.60] at ( 21.60, 58.56) {5};

\node[text=drawColor,anchor=base east,inner sep=0pt, outer sep=0pt, scale=  0.60] at ( 21.60, 90.39) {10};

\node[text=drawColor,anchor=base east,inner sep=0pt, outer sep=0pt, scale=  0.60] at ( 21.60,122.21) {15};

\node[text=drawColor,anchor=base east,inner sep=0pt, outer sep=0pt, scale=  0.60] at ( 21.60,154.04) {20};

\node[text=drawColor,anchor=base east,inner sep=0pt, outer sep=0pt, scale=  0.60] at ( 21.60,185.86) {25};
\end{scope}
\begin{scope}
\path[clip] (  0.00,  0.00) rectangle (332.44,195.13);
\definecolor[named]{drawColor}{rgb}{0.00,0.00,0.00}

\path[draw=drawColor,line width= 0.4pt,line join=round,line cap=round] (258.70,181.56) rectangle (323.00,138.36);
\definecolor[named]{fillColor}{rgb}{1.00,0.00,0.00}

\path[draw=drawColor,line width= 0.4pt,line join=round,line cap=round,fill=fillColor] (264.10,176.16) rectangle (268.42,172.56);
\definecolor[named]{fillColor}{rgb}{0.80,1.00,0.00}

\path[draw=drawColor,line width= 0.4pt,line join=round,line cap=round,fill=fillColor] (264.10,168.96) rectangle (268.42,165.36);
\definecolor[named]{fillColor}{rgb}{0.00,1.00,0.40}

\path[draw=drawColor,line width= 0.4pt,line join=round,line cap=round,fill=fillColor] (264.10,161.76) rectangle (268.42,158.16);
\definecolor[named]{fillColor}{rgb}{0.00,0.40,1.00}

\path[draw=drawColor,line width= 0.4pt,line join=round,line cap=round,fill=fillColor] (264.10,154.56) rectangle (268.42,150.96);
\definecolor[named]{fillColor}{rgb}{0.80,0.00,1.00}

\path[draw=drawColor,line width= 0.4pt,line join=round,line cap=round,fill=fillColor] (264.10,147.36) rectangle (268.42,143.76);

\node[text=drawColor,anchor=base west,inner sep=0pt, outer sep=0pt, scale=  0.60] at (273.82,172.30) {Strongly Agree};

\node[text=drawColor,anchor=base west,inner sep=0pt, outer sep=0pt, scale=  0.60] at (273.82,165.10) {Agree};

\node[text=drawColor,anchor=base west,inner sep=0pt, outer sep=0pt, scale=  0.60] at (273.82,157.90) {Neutral};

\node[text=drawColor,anchor=base west,inner sep=0pt, outer sep=0pt, scale=  0.60] at (273.82,150.70) {Disagree};

\node[text=drawColor,anchor=base west,inner sep=0pt, outer sep=0pt, scale=  0.60] at (273.82,143.50) {Strongly Disagree};
\end{scope}
\end{tikzpicture}

\end{frame}


\begin{frame}[fragile]

\begin{itemize}
\item Participants general comments
\end{itemize}

\frametitle{Curator Interface UX Experiment (3)}
"\textbf{$\cdots$Also, I fail to see how Bonolo differentiates itself from something like Dropbox. I can create a folder structure on my PC and upload it to Dropbox very easily. I can then browse my files and folders in Dropbox's web interface.$\cdots$}"\\

\bigskip

"\textbf{$\cdots$I have to say though that I managed to complete the tasks without watching the video (which is a great sign I think). I'm impatient with manuals but even worse with instructional videos$\cdots$"}
\end{frame}


\begin{frame}[fragile]
\frametitle{Repository Performance Experiment}
\begin{itemize}
\item Objective
\begin{itemize}
\item Impact of file store structure on performance
\item Performance metrics: response time
\end{itemize}
\item Test Environment
\begin{itemize}
\item Intel Core 2 Duo CPU E7400@ 2.80GHz
\item 2 GB RAM
\item 32-bit Windows 7 Ultimate edition
\end{itemize} 
\item Approach
\begin{itemize}
\item Structured and unstructured collections
\item Exponential increase of files in collections
\item Load time and corresponding data transfer during navigation
\end{itemize}
\end{itemize}
\end{frame}


\begin{frame}[fragile]
\frametitle{Repository Performance Experiment (2)}
\centering
% Created by tikzDevice version 0.6.2-92-0ad2792 on 2012-11-10 03:17:14
% !TEX encoding = UTF-8 Unicode
\begin{tikzpicture}[x=1pt,y=1pt]
\definecolor[named]{fillColor}{rgb}{1.00,1.00,1.00}
\path[use as bounding box,fill=fillColor,fill opacity=0.00] (0,0) rectangle (332.44,195.13);
\begin{scope}
\path[clip] (  0.00,  0.00) rectangle (332.44,195.13);
\definecolor[named]{drawColor}{rgb}{0.00,0.00,1.00}

\path[draw=drawColor,line width= 0.4pt,line join=round,line cap=round] ( 38.71, 37.11) --
	(100.67, 36.87) --
	(162.62, 36.99) --
	(224.58, 37.20) --
	(286.53, 37.17);

\path[draw=drawColor,line width= 0.4pt,line join=round,line cap=round] ( 38.71, 37.11) circle (  1.35);

\path[draw=drawColor,line width= 0.4pt,line join=round,line cap=round] (100.67, 36.87) circle (  1.35);

\path[draw=drawColor,line width= 0.4pt,line join=round,line cap=round] (162.62, 36.99) circle (  1.35);

\path[draw=drawColor,line width= 0.4pt,line join=round,line cap=round] (224.58, 37.20) circle (  1.35);

\path[draw=drawColor,line width= 0.4pt,line join=round,line cap=round] (286.53, 37.17) circle (  1.35);
\end{scope}
\begin{scope}
\path[clip] (  0.00,  0.00) rectangle (332.44,195.13);
\definecolor[named]{drawColor}{rgb}{0.00,0.00,0.00}

\path[draw=drawColor,line width= 0.4pt,line join=round,line cap=round] ( 38.71, 28.80) -- (286.53, 28.80);

\path[draw=drawColor,line width= 0.4pt,line join=round,line cap=round] ( 38.71, 28.80) -- ( 38.71, 25.20);

\path[draw=drawColor,line width= 0.4pt,line join=round,line cap=round] (100.67, 28.80) -- (100.67, 25.20);

\path[draw=drawColor,line width= 0.4pt,line join=round,line cap=round] (162.62, 28.80) -- (162.62, 25.20);

\path[draw=drawColor,line width= 0.4pt,line join=round,line cap=round] (224.58, 28.80) -- (224.58, 25.20);

\path[draw=drawColor,line width= 0.4pt,line join=round,line cap=round] (286.53, 28.80) -- (286.53, 25.20);

\node[text=drawColor,anchor=base,inner sep=0pt, outer sep=0pt, scale=  0.60] at ( 38.71, 15.84) {1,024};

\node[text=drawColor,anchor=base,inner sep=0pt, outer sep=0pt, scale=  0.60] at (100.67, 15.84) {2,048};

\node[text=drawColor,anchor=base,inner sep=0pt, outer sep=0pt, scale=  0.60] at (162.62, 15.84) {4,096};

\node[text=drawColor,anchor=base,inner sep=0pt, outer sep=0pt, scale=  0.60] at (224.58, 15.84) {8,192};

\node[text=drawColor,anchor=base,inner sep=0pt, outer sep=0pt, scale=  0.60] at (286.53, 15.84) {16,384};

\path[draw=drawColor,line width= 0.4pt,line join=round,line cap=round] ( 28.80, 34.69) -- ( 28.80,182.04);

\path[draw=drawColor,line width= 0.4pt,line join=round,line cap=round] ( 28.80, 34.69) -- ( 25.20, 34.69);

\path[draw=drawColor,line width= 0.4pt,line join=round,line cap=round] ( 28.80, 64.16) -- ( 25.20, 64.16);

\path[draw=drawColor,line width= 0.4pt,line join=round,line cap=round] ( 28.80, 93.63) -- ( 25.20, 93.63);

\path[draw=drawColor,line width= 0.4pt,line join=round,line cap=round] ( 28.80,123.10) -- ( 25.20,123.10);

\path[draw=drawColor,line width= 0.4pt,line join=round,line cap=round] ( 28.80,152.57) -- ( 25.20,152.57);

\path[draw=drawColor,line width= 0.4pt,line join=round,line cap=round] ( 28.80,182.04) -- ( 25.20,182.04);

\node[text=drawColor,anchor=base east,inner sep=0pt, outer sep=0pt, scale=  0.60] at ( 21.60, 32.63) {0};

\node[text=drawColor,anchor=base east,inner sep=0pt, outer sep=0pt, scale=  0.60] at ( 21.60, 62.10) {1000};

\node[text=drawColor,anchor=base east,inner sep=0pt, outer sep=0pt, scale=  0.60] at ( 21.60, 91.56) {2000};

\node[text=drawColor,anchor=base east,inner sep=0pt, outer sep=0pt, scale=  0.60] at ( 21.60,121.03) {3000};

\node[text=drawColor,anchor=base east,inner sep=0pt, outer sep=0pt, scale=  0.60] at ( 21.60,150.50) {4000};

\node[text=drawColor,anchor=base east,inner sep=0pt, outer sep=0pt, scale=  0.60] at ( 21.60,179.97) {5000};
\end{scope}
\begin{scope}
\path[clip] (  0.00,  0.00) rectangle (332.44,195.13);
\definecolor[named]{drawColor}{rgb}{1.00,0.00,0.00}

\path[draw=drawColor,line width= 0.4pt,dash pattern=on 4pt off 4pt ,line join=round,line cap=round] ( 38.71, 44.27) --
	(100.67, 52.20) --
	(162.62, 70.35) --
	(224.58, 94.51) --
	(286.53,158.46);

\path[draw=drawColor,line width= 0.4pt,line join=round,line cap=round] ( 37.52, 43.07) rectangle ( 39.91, 45.47);

\path[draw=drawColor,line width= 0.4pt,line join=round,line cap=round] ( 99.47, 51.00) rectangle (101.86, 53.39);

\path[draw=drawColor,line width= 0.4pt,line join=round,line cap=round] (161.42, 69.15) rectangle (163.82, 71.55);

\path[draw=drawColor,line width= 0.4pt,line join=round,line cap=round] (223.38, 93.32) rectangle (225.77, 95.71);

\path[draw=drawColor,line width= 0.4pt,line join=round,line cap=round] (285.33,157.26) rectangle (287.73,159.66);
\definecolor[named]{drawColor}{rgb}{0.00,0.00,0.00}

\node[text=drawColor,anchor=base,inner sep=0pt, outer sep=0pt, scale=  0.60] at (162.62,  1.44) {Files in Directory};

\node[text=drawColor,rotate= 90.00,anchor=base,inner sep=0pt, outer sep=0pt, scale=  0.60] at (  5.76,108.36) {Time (Milliseconds)};

\path[draw=drawColor,line width= 0.4pt,line join=round,line cap=round] ( 38.71,182.04) rectangle ( 94.12,160.44);
\definecolor[named]{drawColor}{rgb}{0.00,0.00,1.00}

\path[draw=drawColor,line width= 0.4pt,line join=round,line cap=round] ( 40.33,174.84) -- ( 51.13,174.84);
\definecolor[named]{drawColor}{rgb}{1.00,0.00,0.00}

\path[draw=drawColor,line width= 0.4pt,dash pattern=on 4pt off 4pt ,line join=round,line cap=round] ( 40.33,167.64) -- ( 51.13,167.64);
\definecolor[named]{drawColor}{rgb}{0.00,0.00,1.00}

\path[draw=drawColor,line width= 0.4pt,line join=round,line cap=round] ( 45.73,174.84) circle (  1.35);
\definecolor[named]{drawColor}{rgb}{1.00,0.00,0.00}

\path[draw=drawColor,line width= 0.4pt,line join=round,line cap=round] ( 44.54,166.44) rectangle ( 46.93,168.83);
\definecolor[named]{drawColor}{rgb}{0.00,0.00,0.00}

\node[text=drawColor,anchor=base west,inner sep=0pt, outer sep=0pt, scale=  0.60] at ( 56.53,172.77) {Structured};

\node[text=drawColor,anchor=base west,inner sep=0pt, outer sep=0pt, scale=  0.60] at ( 56.53,165.57) {Unstructured};
\end{scope}
\end{tikzpicture}

\end{frame}

\begin{frame}[fragile]
\frametitle{Conclusion}

\begin{itemize}
\item Experimental results look promising
\begin{itemize}
\item Effectiveness
\item Usability
\item Medium-sized collections
\end{itemize}

\item Work in progress
\begin{itemize}
\item Evaluation
\begin{itemize}
\item Flexibility
\item Scalability
\end{itemize}
\end{itemize}

\item Future work
\begin{itemize}
\item Reference implementation
\begin{itemize} 
\item Design principles
\item Extensibility
\end{itemize}
\end{itemize}

\end{itemize}

\end{frame}

\begin{frame}[plain]
%\frametitle{Thank You}
\begin{center}
\bigskip
\textbf{\fontsize{18}{18}\selectfont Thank You} \\
\vspace{\stretch{0.5}}
\textbf{\fontsize{18}{18}\selectfont Questions?} \\
\vspace{\stretch{0.5}}
\textbf{\fontsize{18}{18}\selectfont Additional Information} \\
\bigskip
\textbf{\fontsize{18}{18}\selectfont \url{http://dl.cs.uct.ac.za}}
\end{center}
\end{frame}

\end{document}
